% =====================================================================
% 第 1 节  引言  （原书 p.1；对应 work/pages/page_001.png）
% =====================================================================

\section{引言}\label{sec:1}

\begin{keytakeaways}
  \item \textbf{物理图像}：量子场论有三种等价而互补的表述——拉格朗日/路径积分、哈密顿量、算符代数；其中算符代数进路最内蕴，而共形场论（CFT）正是它最清晰的用武之地。
  \item \textbf{数学工具}：态-算符对应、模不变性、重整化群流；二维情形还多出无限维共形对称群。
  \item \textbf{本节结论}：大量庞加莱不变的局域量子场论都可视为从紫外 CFT 出发的 RG 流，因此 CFT 的分类与"求解"是量子场论的核心问题；二维 CFT 的版图远比 Calabi--Yau 模型所覆盖的角落辽阔。
\end{keytakeaways}

量子场论是现代理论物理学的支柱，其基础却建立得相当不能令人满意。本质上存在三条进路：（1）拉格朗日量与路径积分；（2）带紫外正规化的哈密顿量；（3）算符代数。拉格朗日进路适用于构造有效理论、发展微扰论，有时也能超越微扰论来处理受对称性保护的观测量，并可通过欧氏时空的格点离散化给出非微扰表述。哈密顿量进路在实用上或借助空间格点离散化、或借助截断共形空间方法而发挥作用。算符代数进路是量子场论最内蕴的表述，而它在共形场论（CFT）的局域算符语境中得到最好的理解。很大一类庞加莱不变的局域量子场论可以描述为从紫外 CFT 出发的重整化群流。\footnote{有一种流传的说法称：所有量子场论都来自紫外处的 CFT；这是不对的——标准反例是二、三维的规范场论。}因此，CFT 的分类与"求解"对量子场论的研究十分重要。\footnote{所谓"求解"一个 CFT，通常是指确定其局域算符谱与算符乘积展开代数。}

二维共形场论享有无限维的共形对称群。它们还受到模不变性这一重要约束，后者使二维 CFT 在任意曲面上都有自洽的定义。尽管已知大量例子，我们距离分类或求解大多数二维 CFT 仍很遥远。例如，每个 Calabi--Yau 流形都通过非线性 sigma 模型给出一个超共形场论，其算符谱在精确边缘形变模空间的一般点上，除受超对称保护的部分外通常并不知晓；此外，Calabi--Yau 模型也只占据二维 CFT 版图的一个极小角落。

在本讲义中，我将勾勒二维 CFT 的算符定义，讨论若干标准例子，并概述可用于分类与求解 CFT 的自举方法。

% [Tutor Note: 把 CFT 的"解"归结为"谱 + OPE 系数"两样数据，是自举（bootstrap）纲领的出发点：一旦给定全部原初算符的标度维数与三点结构常数，交叉对称性与模不变性就变成对这些数据的代数/泛函方程，而不再需要路径积分。本讲义第 5 节正是沿这条思路，把约束组织成半正定规划来数值求解。读完全文后可回头对照此处的"三条进路"，会更容易看清算符代数进路为何是自举的自然语言。]
